% !PRINTED 30
\weberpage{30}
\begin{center}
\S\,8. \textbf{\gesperrt{Das Jacobische Transformationsprinzip.}}
\end{center}

Die \gesperrt{Jacobische} Transformationstheorie stellt sich im allgemeinen die
Aufgabe, ein elliptisches Differential durch eine algebraische
Substitution auf ein anderes elliptisches Differential zurückzuführen.
Indem wir hier die Grundlagen dieser Theorie auseinandersetzen, bedienen
wir uns der Darstellung des elliptischen Differentials in homogenen
Variablen [\S\,1, (4)]:
\weq{1}{\frac{x\,dy - y\,dx}{\sqrt{f(x,y)}}\,,}
und substituieren darin für $x,y$ zwei teilerfremde ganze homogene
Funktionen gleichen, aber beliebigen Grades $n$ zweier neuer Variablen
$\xi,\eta$:
\weq{2}{x = U(\xi,\eta), \qquad y = V(\xi,\eta).}
Aus den Gleichungen
\[
\begin{aligned}
nU &= \xi\,\frac{\partial U}{\partial \xi} + \eta\,\frac{\partial U}{\partial \eta}, &
dU &= d\xi\,\frac{\partial U}{\partial \xi} + d\eta\,\frac{\partial U}{\partial \eta},\\
nV &= \xi\,\frac{\partial V}{\partial \xi} + \eta\,\frac{\partial V}{\partial \eta}, &
dV &= d\xi\,\frac{\partial V}{\partial \xi} + d\xi\,\frac{\partial V}{\partial \eta} % [sic] second differential printed d\xi; the total differential requires d\eta
\end{aligned}
\]
folgt sodann
\weq{3}{U\,dV - V\,dU = H(\xi\,d\eta - \eta\,d\xi),}
wenn $H$ die Funktionaldeterminante
\weq{4}{H = \frac{1}{n}\left(\frac{\partial U}{\partial \xi}\frac{\partial V}{\partial \eta} - \frac{\partial V}{\partial \xi}\frac{\partial U}{\partial \eta}\right)}
bedeutet.

Danach geht das elliptische Differential (1) in das folgende über:
\weq{5}{H\,\frac{\xi\,d\eta - \eta\,d\xi}{\sqrt{f(U,V)}}\,.}
Damit nun dieses Differential wieder die Form eines elliptischen erhält,
ist erforderlich, daß von der Funktion $f(U,V)$, deren Grad der $4n$te
ist, sich ein quadratischer Faktor $T^2$ vom Grade $4n-4$ absondern
lasse, also, wenn $\varphi(\xi,\eta)$ eine Funktion vierten Grades
bedeutet, daß
\weq{6}{f(U,V) = T^2\,\varphi(\xi,\eta)}
werde, wodurch, wenn
\weq{7}{\frac{T}{H} = M}
%CONT>>
